% SPDX-License-Identifier: Apache-2.0
% covers: Pwm
\datasheet{Pwm}{Four pulse width modulated outputs on AXI-Lite: one counter, a duty a channel, and a shadow so a width never changes under a pulse.}

\dsfacts{\code{lib/parts/src/pwm.rs}}{\code{txhdl\_parts::pwm::Pwm}}%
{\code{//docs:examples}}

\subsection*{Function}

A pin is on or off, and much of what a board drives wants a value
between: an LED at a third of its brightness, a fan at half its speed, a
servo told where to stand. \code{Pwm} switches four pins fast enough
that the load cannot follow, and what the load sees is the fraction of
each period a pin was high.

One counter runs the period for all four channels, and a channel is
high while the counter is below its duty. That comparison is what makes
the ends flat: a duty of zero never wins it and a duty of the period or
more always does, so neither end leaves a pulse one cycle wide.

\subsection*{Parameters}

None. The period, the duties and the four polarity bits are registers,
so one instance drives four channels of different widths; the channel
count and the 16-bit period are fixed in the unit.

\subsection*{Register map}

Offsets from the peripheral's base. The part selects the word by bits 4
to 2 of the address and looks at no other bits, so the bridge in front
of it decides the base and the map repeats every 32 bytes.
\code{pwm::duty(i)} gives channel \code{i}'s offset, which is
\code{8 + 4i}.

\begin{center}\footnotesize
\begin{tabular}{@{}llp{0.5\textwidth}@{}}
\toprule
Offset & Word & Access \\
\midrule
\code{0x00} & \code{ctrl} & Read and write, eight bits. Bit 0 is
\code{CTRL\_ENABLE}, bit 1 \code{CTRL\_CENTRE}, and bits 4 to 7 the
polarity of channels 0 to 3, which \code{pwm::polarity(i)} names. \\
\code{0x04} & \code{period} & Read and write. A read returns the
shadow, not the period the counter is running. \\
\code{0x08} & \code{duty0} & Read and write, the shadow again. \\
\code{0x0c} & \code{duty1} & Read and write. \\
\code{0x10} & \code{duty2} & Read and write. \\
\code{0x14} & \code{duty3} & Read and write. \\
\code{0x18} & \code{count} & Read only, and it is the live counter
rather than a shadow. A write is taken and does nothing. \\
\code{0x1c} & & Reads as zero. \\
\bottomrule
\end{tabular}
\end{center}

A write takes the low 16 bits of the word for a period or a duty and
the low 8 for \code{ctrl}, and the rest are dropped.

\dstables{Pwm}

\subsection*{Behaviour}

\begin{itemize}
\item Edge aligned, which is \code{CTRL\_CENTRE} clear, the counter
runs from zero to the period less one and returns to zero, so a period
is that many cycles and every pulse starts at the same moment.
\item Centre aligned it counts up and then down, holding one cycle at
each end, so a period is twice the count and a pulse is centred in it,
as one run rather than two at the ends. A bank of outputs driven that
way does not switch all at once, which is what a supply prefers.
\item A write of a period or a duty goes to a shadow, and every shadow
is loaded together when the counter wraps. A width that changed under
the counter would make one pulse of neither the old width nor the new,
which a load that averages sees as a step nobody asked for and a servo
reads as a command.
\item While the counter is stopped the shadows are loaded every cycle
and the counter is held at zero, so a period written then takes effect
at once: no pulse is under way to spoil.
\item A channel's pin is the comparison exclusive-or'd with its
polarity bit. While the counter is stopped the pins are the four
polarity bits themselves, so a design chooses what a stopped output
rests at.
\item A period of zero stops the counter, since a period shorter than a
cycle is no period at all.
\item A read is answered on \code{bus\_r} in the cycle its address beat is
taken, with \code{OKAY}. A write needs its address beat, its data beat
and room on \code{bus\_b} in the same cycle, and is answered \code{OKAY} in
that cycle.
\end{itemize}

\subsection*{Verification}

\begin{itemize}
\item Six tests in \code{lib/parts/src/pwm.rs}, in
\code{//lib/parts:parts\_test}, drive it over AXI-Lite and count the
cycles each pin was high.
\code{a\_duty\_is\_the\_cycles\_high\_in\_a\_period} finds three and seven
of ten, four periods over.
\code{the\_ends\_are\_flat} finds a duty of zero never high and a duty of
the period, or past it, never low.
\code{polarity\_turns\_a\_channel\_over} finds the same duty high for
nine cycles and for twenty-one.
\code{a\_duty\_written\_mid\_period\_makes\_no\_runt} widens a duty in the
middle of a period and finds every whole pulse two wide or eight, and
the last of them eight.
\code{a\_centred\_pulse\_is\_one\_run\_a\_period} finds four of eight,
centred, as one run of eight in sixteen.
\code{a\_stopped\_counter\_sits\_at\_the\_polarity} finds the pins at
their polarity bits with the enable clear.
\item \code{//lib/examples:ex\_pwm} drives it through a
\code{LiteBridge1} on an AXI4 link and prints what the program
measured: the duties as cycles high, the pulses around a change of
width, a channel turned over, and the centred pulses.
\item The build simulates the netlist against that run under nvc and
Verilator: \code{//docs:pwm\_sim\_pwm\_tb\_test} and
\code{//docs:pwm\_vsim\_test}.
\end{itemize}

\subsection*{Used by}

Nothing in the tree yet. Issue 248 asks for a program on the board that
fades an LED with it.

\subsection*{Limits}

Four channels of 16 bits, neither a parameter. The period is in system
clock cycles, since there is no prescaler, so the slowest period is
65\,535 cycles. Byte strobes are ignored. Nothing is refused: an
address that names no word reads as zero or as the counter, and every
read and write is answered \code{OKAY}. A read of a period or a duty
returns what was written rather than what the counter is using, so
there is no way to read the value in use. There is no interrupt, no
dead time between complementary outputs, no fault input, no phase
between channels and no period of its own for any of them. A period of
zero written while the counter is enabled cannot take effect, because
the load of the shadows is itself gated on the counter running; the
enable has to come off first.
